% =====================================================================
% CHAPTER 1: INTRODUCTION
% =====================================================================
\chapter{Introduction}

\section{Overview}
The rapid advancement of Artificial Intelligence and embedded systems has opened new frontiers in wearable technology. Smart glasses represent one of the most promising applications at the intersection of these technologies, offering the potential to augment human perception through real-time computational assistance. This project will present the design and implementation of AI-powered smart glasses that leverage the ESP32-CAM microcontroller as the primary hardware platform, combined with a hybrid cloud-local AI backend architecture to deliver sophisticated visual and linguistic assistance capabilities.

The ESP32-CAM is a compact, low-cost development board featuring a dual-core Xtensa LX6 processor running at 240 MHz with 520 KB of SRAM and 4 MB of flash memory. The module includes an integrated OV2640 camera capable of capturing images at resolutions up to 1600$\times$1200 pixels. For this project, images will be captured at 640$\times$480 (VGA) resolution to balance quality with transmission speed over wireless networks. The ESP32-CAM supports IEEE 802.11 b/g/n Wi-Fi and Bluetooth 4.2, enabling seamless wireless communication with the backend server running on a local machine.

The backend server will implement a novel Smart Routing architecture that intelligently distributes incoming AI tasks across multiple processing pathways. Computationally intensive tasks such as image understanding and natural language question answering will be offloaded to cloud-based large language models accessed through the Groq API free tier. This tier provides access to state-of-the-art models including Llama 4 Scout 17B for multimodal vision tasks and Llama 3.3 70B Versatile for language processing, with a generous allowance of 14,400 free API calls per day. These cloud-based models will enable the system to perform advanced reasoning, scene understanding, and multilingual translation without imposing any computational burden on the local hardware.

Lightweight tasks that benefit from low latency and offline capability will be processed locally using open-source AI models. EasyOCR will handle optical character recognition for extracting text from images, supporting over 80 languages with robust accuracy. OpenAI Whisper in its tiny configuration (39 million parameters) will provide speech-to-text transcription capability with minimal resource requirements. The Smart Router will ensure that only one local model occupies memory at any given time, and models will be explicitly unloaded after each request. This memory management strategy is expected to keep peak RAM usage under 600 MB, making the system compatible with entry-level NVIDIA GPUs such as the GTX 1650 and RTX 3050.

The integration of multiple AI modalities within a single wearable device represents a significant step forward in affordable assistive technology. By combining computer vision, natural language processing, optical character recognition, and automatic speech recognition, the system will provide comprehensive assistance that goes far beyond what traditional assistive devices such as white canes or basic electronic travel aids can offer.

\section{Purpose}
The primary purpose of this project is to develop an affordable, functional smart glasses prototype that demonstrates the practical application of AI-powered assistive technology. The project will aim to bridge the gap between expensive commercial smart glasses and the needs of users who require visual assistance, particularly visually impaired individuals in developing countries where cost is a significant barrier to adoption.

Commercial smart glasses such as Google Glass Enterprise Edition and Meta Ray-Ban Smart Glasses are priced well beyond the reach of most users in India and other developing nations. By utilizing the ESP32-CAM platform and free-tier cloud AI services, this project will demonstrate that comparable AI capabilities can be delivered at a fraction of the cost using readily available components. The system will provide five core functionalities---scene description, text extraction, speech recognition, question answering, and translation---all accessible through a simple button-based interface mounted on a standard glasses frame.

The project will also serve as a proof-of-concept for the hybrid cloud-local AI architecture, demonstrating that sophisticated AI workloads can be efficiently distributed between cloud and edge resources without requiring expensive dedicated hardware or paid cloud subscriptions. This architectural pattern will have broader implications for IoT device design, as it enables resource-constrained edge devices to access powerful AI capabilities while maintaining low latency for time-critical operations.

Furthermore, the project will contribute to the field of accessible technology by providing an open-source, replicable design that can be adopted and customized by researchers, educators, and humanitarian organizations working to improve the lives of visually impaired individuals. The use of open-source software components and widely available hardware ensures that the design will be accessible to a broad audience.

\section{Project Relevance}
This project is relevant to multiple domains within computer science and engineering, reflecting the interdisciplinary nature of modern IoT and AI applications:

\begin{itemize}[leftmargin=1.5cm]
\item \textbf{Internet of Things (IoT):} The system will demonstrate core IoT principles through wireless sensor data collection (camera, microphone), cloud processing via RESTful APIs, and actuator feedback through the OLED display and speaker. The ESP32-CAM will serve as an IoT edge node that communicates with a local gateway server.

\item \textbf{Embedded Systems:} The ESP32-CAM firmware development will involve real-time programming, peripheral management (camera, I2C for OLED, I2S for microphone), interrupt-driven button handling, and resource-constrained optimization. The firmware will need to manage multiple communication protocols simultaneously while maintaining responsiveness.

\item \textbf{Artificial Intelligence:} The project will integrate multiple AI paradigms including computer vision (scene understanding via multimodal LLMs), natural language processing (question answering and translation), optical character recognition (text extraction from images), and automatic speech recognition (voice-to-text transcription). The Smart Routing architecture will demonstrate practical AI model lifecycle management.

\item \textbf{Cloud Computing:} The hybrid architecture will demonstrate practical cloud-edge computing patterns, RESTful API integration, and intelligent workload distribution strategies. The use of the Groq API free tier will showcase how cloud AI services can be leveraged cost-effectively for resource-constrained applications.

\item \textbf{Assistive Technology:} The system will directly address the needs of visually impaired users, aligning with the United Nations Sustainable Development Goal 10 (Reduced Inequalities) and Goal 3 (Good Health and Well-being). The dual output modality (visual display and audio speaker) will ensure accessibility for users with varying degrees of visual impairment.

\item \textbf{Cybersecurity and Privacy:} The edge-processing approach for sensitive tasks (OCR, speech) will keep personal data local, while cloud API calls will transmit only encoded image data over HTTPS-encrypted channels. This privacy-aware architecture will demonstrate responsible AI deployment practices in wearable devices.
\end{itemize}
